\chapter{数学与概率预备知识}

\begin{quote}\small
\textit{本章摘要。}本章建立全书共同使用的数学语言。先从向量、矩阵和微积分说明模型如何表示对象、变化与优化，再引入概率、分布和期望，最后解释熵、距离、降维和数值稳定性。阅读顺序对应后文的依赖关系：先能读懂确定性模型，再能处理随机性，最后才能理解贝叶斯推断、深度学习和流形表示中的公式。
\end{quote}

本书从数学开始，并不是要求读者先成为数学家，而是要建立一种可核查的阅读习惯：每个运算必须作用在合适的对象上，每个结论必须有成立条件。设备的一组传感器读数先成为向量，多个样本组成矩阵；预测误差成为标量函数，其变化由梯度描述；同一工况下仍会出现不同结果，则需要条件概率。以下工具沿着这条路径出现，而不是彼此孤立的公式表。

可以始终带着同一组设备读数来理解这些工具。把温度和振动画成一个点，是几何表示；调整权重让预测更接近测量值，需要导数；重复运行后读数仍然有波动，就要用概率和期望。先看清每一步输入和输出的维度，再跟着推导计算。后面遇到具体模型时，可以回到这里核对它用了哪些条件。

\section{为什么需要数学语言}
机器学习不是把公式堆在代码旁边，而是用一套精确语言说明四件事：对象是什么，假设是什么，优化什么，以及结果在哪些条件下成立。本章只介绍后文反复使用的工具。读者不必一次记住所有细节，但应能在看到一个公式时识别其输入、输出、随机性与求和范围。

\section{向量、矩阵与几何}
向量$\bm{x}\in\mathbb{R}^d$可以表示一个样本，也可以表示参数、梯度或隐状态。内积$\bm{x}^{\mathsf T}\bm{y}$衡量方向相似性，范数$\|\bm{x}\|_2=(\sum_jx_j^2)^{1/2}$衡量长度。矩阵$A\in\mathbb{R}^{m\times d}$是从$d$维空间到$m$维空间的线性变换。其列空间、零空间和秩分别描述可表达方向、被压缩方向与有效维数。

对称半正定矩阵满足$\bm{v}^{\mathsf T}A\bm{v}\geq0$。协方差矩阵必须具有这一性质，因为方差不可能为负。二次型$\bm{v}^{\mathsf T}A\bm{v}$还可以看作沿方向$\bm{v}$观察曲率，这使 Hessian 矩阵成为优化和不确定性分析的核心对象。

\paragraph{维度、投影与半正定性的由来。}
若$A\in\R^{m\times d}$、$\bm x\in\R^d$，则$(A\bm x)_i=\sum_{j=1}^d A_{ij}x_j$，输出有$m$个分量。内积不是不受尺度影响的相似度：将一个向量放大两倍会使内积也放大两倍；只有先归一化，内积才等于夹角余弦。对单位向量$\bm u$，把$\bm x$投影到其张成的直线上，需要最小化
\begin{align}
 \|\bm x-a\bm u\|_2^2
 &=\|\bm x\|_2^2-2a\bm u^{\mathsf T}\bm x+a^2,\\
 \frac{d}{da}\|\bm x-a\bm u\|_2^2&=-2\bm u^{\mathsf T}\bm x+2a=0.\label{eq:math-projection-stationarity}
\end{align}
由式\eqref{eq:math-projection-stationarity}得$a=\bm u^{\mathsf T}\bm x$，投影为$\bm u\bm u^{\mathsf T}\bm x$，残差与$\bm u$正交。若$\bm u$不是单位向量，系数必须再除以$\bm u^{\mathsf T}\bm u$。

设随机向量$\bm X$有有限二阶矩，$\bm\mu=\E\bm X$，定义$\Sigma=\E[(\bm X-\bm\mu)(\bm X-\bm\mu)^{\mathsf T}]\in\R^{d\times d}$。对任意$\bm v\in\R^d$，
\begin{equation}
 \bm v^{\mathsf T}\Sigma\bm v
 =\E\bigl[(\bm v^{\mathsf T}(\bm X-\bm\mu))^2\bigr]
 =\Var(\bm v^{\mathsf T}\bm X)\geq0.\label{eq:math-covariance-psd}
\end{equation}
式\eqref{eq:math-covariance-psd}同时证明了协方差半正定，并说明二次型就是投影后的方差。半正定并不保证可逆：若某个传感器恰好是另一个的两倍，就存在非零方向上的方差为零。

\section{微积分与链式法则}
若标量函数$J(\bm{\theta})$对参数可微，梯度为
\begin{equation}
 \nabla_{\bm{\theta}}J=\left[\frac{\partial J}{\partial\theta_1},\ldots,\frac{\partial J}{\partial\theta_d}\right]^{\mathsf T}.
\end{equation}
梯度指向函数增长最快的方向，因此$-\nabla J$是局部下降方向。对复合函数$J(f(\bm{\theta}))$，链式法则给出
\begin{equation}
 \frac{\partial J}{\partial\theta_j}=\sum_i\frac{\partial J}{\partial f_i}\frac{\partial f_i}{\partial\theta_j}.
\end{equation}
反向传播正是以高效的图结构方式反复应用链式法则，而不是一种神秘的独立学习规则。

\paragraph{为什么负梯度是局部下降方向。}
令$\bm g=\nabla J(\bm\theta)$。可微性给出$J(\bm\theta+\epsilon\bm v)=J(\bm\theta)+\epsilon\bm g^{\mathsf T}\bm v+o(\epsilon)$。在$\|\bm v\|_2=1$的方向中，Cauchy--Schwarz不等式给出$\bm g^{\mathsf T}\bm v\geq-\|\bm g\|_2$，当$\bm g\ne0$时由$\bm v=-\bm g/\|\bm g\|_2$取到等号。这只说明足够小的步长能下降，不保证任意步长都有效，也不保证梯度为零时已经达到最小值。

以$\bm r=A\bm\theta-\bm y\in\R^m$和$J=\tfrac12\bm r^{\mathsf T}\bm r$为例，微分逐步给出
\begin{equation}
 dJ=\bm r^{\mathsf T}d\bm r
 =\bm r^{\mathsf T}A\,d\bm\theta
 =(A^{\mathsf T}\bm r)^{\mathsf T}d\bm\theta,
 \qquad \nabla J=A^{\mathsf T}(A\bm\theta-\bm y)\in\R^d.
\end{equation}
一般地，若$\bm f:\R^d\to\R^m$的Jacobian为$D\bm f\in\R^{m\times d}$，则$\nabla_{\bm\theta}J=(D\bm f)^{\mathsf T}\nabla_{\bm f}J$。转置把输出空间的误差信号拉回参数空间，维度检查可以帮助发现漏转置和广播错误。

\section{概率、条件概率与期望}
随机变量$X$的分布描述不确定性。联合分布满足乘法法则
\begin{equation}
 p(x,y)=p(y\mid x)p(x),
 \qquad p(y\mid x)=\frac{p(x,y)}{p(x)}.\label{eq:math-conditional-probability}
\end{equation}
边缘化把未知变量积分掉：$p(x)=\int p(x,y)\,dy$。期望是随机量的加权平均，方差是
\begin{equation}
 \Var(X)=\mathbb{E}[(X-\mathbb{E}[X])^2].
\end{equation}
条件期望$\mathbb{E}[Y\mid X]$是给定输入后对输出的最佳平方损失预测，这解释了为什么均方误差回归在理想条件下学习条件均值。

\paragraph{条件均值为何最优。}
对离散变量使用概率质量，对连续变量使用密度；式\eqref{eq:math-conditional-probability}的条件概率比值只在分母为正处使用，连续情形不是把单点概率相除。假设$\E[Y^2]<\infty$，固定输入$x$，记$\mu(x)=\E[Y\mid X=x]$。对任意实数预测$a$，把误差拆为$(Y-\mu)+(\mu-a)$，得到
\begin{align}
 \E[(Y-a)^2\mid X=x]
 &=\E[(Y-\mu)^2\mid X=x]
   +2(\mu-a)\E[Y-\mu\mid X=x]+(\mu-a)^2\\
 &=\Var(Y\mid X=x)+(\mu(x)-a)^2.\label{eq:math-conditional-mse}
\end{align}
式\eqref{eq:math-conditional-mse}的第一项与$a$无关，交叉项因条件残差均值为零而消失，所以最优$a$为$\mu(x)$。这不是说有限样本拟合一定得到条件均值，而是说它是无限制平方损失预测的理想目标。

进一步对$X$取期望，并对$Y-\E Y$作同样分解，可得全方差公式
\begin{equation}
 \Var(Y)=\E[\Var(Y\mid X)]+\Var(\E[Y\mid X]).\label{eq:math-total-variance}
\end{equation}
式\eqref{eq:math-total-variance}将总波动分为给定输入后仍存在的波动，以及不同输入导致的均值变化。缺少关键传感器时，前一项可能很大；仅靠更精细地优化参数无法消除这种信息缺失。

\section{高斯分布与最大熵}
多元高斯分布写作
\begin{equation}
 p(\bm{x})=\frac{1}{(2\pi)^{d/2}|\Sigma|^{1/2}}
 \exp\left[-\frac12(\bm{x}-\bm{\mu})^{\mathsf T}\Sigma^{-1}(\bm{x}-\bm{\mu})\right].\label{eq:math-gaussian-density}
\end{equation}
均值决定中心，协方差决定椭球的方向和尺度。给定均值和协方差，高斯分布具有最大的微分熵，因此它常作为“只知道一二阶统计量时不过度添加假设”的基线。

\paragraph{密度公式与最大熵结论的条件。}
式\eqref{eq:math-gaussian-density}的普通密度要求$\Sigma$正定。写$\Sigma=U\Lambda U^{\mathsf T}$，令$\bm z=\Lambda^{-1/2}U^{\mathsf T}(\bm x-\bm\mu)$，则指数中的二次型为$\|\bm z\|_2^2$，变量替换的体积因子为$|\Sigma|^{1/2}$。因此它实际上是独立标准正态密度经旋转、伸缩和平移得到的分布。若$\Sigma$奇异，概率集中在较低维子空间上，不能直接套用含$\Sigma^{-1}$的$d$维密度。

设$q$为具有给定均值和正定协方差的高斯密度，$p$为具有相同矩且相关积分有限的另一密度。因为$\log q$只是常数加二次型，相同的一二阶矩意味着$\E_p\log q=\E_q\log q$。利用下一节证明的Kullback--Leibler divergence (KL)非负性，
\begin{equation}
 0\leq\KL(p\|q)=\int p\log p-\int p\log q
 =-h(p)+h(q),\qquad
 h(q)=\frac12\log\bigl((2\pi e)^d|\Sigma|\bigr).
\end{equation}
故$h(p)\leq h(q)$。最大熵说的是固定矩约束下不再添加额外分布结构，并不证明现实数据必为高斯。微分熵还会随坐标单位改变，不应直接当成离散事件数目。

\section{熵、交叉熵与 KL 散度}
熵的信息度量来自 Shannon 的信息论框架，相对熵的统计区分意义可追溯到 Kullback 与 Leibler 的工作\cite{shannon1948information,kullback1951information}。以下公式区分随机性、编码代价与分布差异，三者不能仅因为都使用对数就被视为同一个量。
离散分布$p$的熵为$H(p)=-\sum_xp(x)\log p(x)$。用$q$编码来自$p$的数据时，平均代价是交叉熵$H(p,q)=-\sum_xp(x)\log q(x)$。两者关系为
\begin{equation}
 H(p,q)=H(p)+\KL(p\|q),
 \qquad \KL(p\|q)=\sum_xp(x)\log\frac{p(x)}{q(x)}\geq0.\label{eq:math-cross-entropy-kl}
\end{equation}
由式\eqref{eq:math-cross-entropy-kl}可知，最小化交叉熵等价于让预测分布接近真实分布。KL 散度不对称，选择$\KL(p\|q)$还是$\KL(q\|p)$会带来不同的近似行为，这一点在变分推断中非常重要。

\begin{example}
 若故障类别的真实分布为$(0.9,0.1)$，模型给出$(0.6,0.4)$，交叉熵会惩罚模型对高概率类别的低估。若模型输出概率而不是只输出类别，系统还可以根据置信度设置不同告警阈值。
\end{example}

\paragraph{KL非负性与交叉熵分解。}
约定$0\log0=0$，若$p(x)>0$但$q(x)=0$，则KL为无穷大。其余情形利用$-\log u\geq1-u$，逐项得到
\begin{equation}
 \KL(p\|q)=\sum_{p(x)>0}p(x)\left[-\log\frac{q(x)}{p(x)}\right]
 \geq\sum_{p(x)>0}(p(x)-q(x))\geq0.
\end{equation}
等号要求两分布相同。再把$\log(p/q)=\log p-\log q$代回求和，立即得到$\KL=H(p,q)-H(p)$。训练时真实分布$p$固定，因而最小化交叉熵与最小化这个方向的KL有相同最优解；用有限样本平均替代期望后，还存在采样误差和模型表达能力限制。

\section{从几何距离到概率距离}
机器学习中常见的“距离”并不只有欧氏距离。欧氏距离比较两个向量在坐标空间中的位置，余弦距离比较方向，Mahalanobis 距离则利用协方差修正不同方向的尺度：
\begin{equation}
 d_M(x,y)=\sqrt{(x-y)^{\mathsf T}\Sigma^{-1}(x-y)}.
\end{equation}
当某个传感器噪声很大时，沿该方向的差异不应与低噪声方向等量齐观。下面列出几种常见距离。它们比较的对象相同，但对“差异”的理解不同：

\begin{align}
 d_1(x,y)&=\|x-y\|_1=\sum_{j=1}^{d}|x_j-y_j|,\label{eq:math-manhattan-distance}\\
 d_2(x,y)&=\|x-y\|_2=\sqrt{\sum_{j=1}^{d}(x_j-y_j)^2},\label{eq:math-euclidean-distance}\\
 d_p(x,y)&=\|x-y\|_p=\left(\sum_{j=1}^{d}|x_j-y_j|^p\right)^{1/p},\quad p\geq1,\label{eq:math-minkowski-distance}\\
 d_\infty(x,y)&=\|x-y\|_\infty=\max_{1\leq j\leq d}|x_j-y_j|.\label{eq:math-chebyshev-distance}
\end{align}

式\eqref{eq:math-manhattan-distance}把各坐标的绝对差直接相加，适合把每个传感器的偏差看成可以分别累积的误差。式\eqref{eq:math-euclidean-distance}对较大的单项差异惩罚更重，常用于连续特征经过同尺度处理后的比较。式\eqref{eq:math-minkowski-distance}把前两者统一起来：$p=1$得到曼哈顿距离，$p=2$得到欧氏距离；当$p$增大时，最大坐标差异的影响逐渐增强。式\eqref{eq:math-chebyshev-distance}只关注最不一致的那个坐标，例如设备的多个安全指标中只要有一项超出容许范围，就把该项作为整体差异的代表。

如果只关心方向而不关心向量长度，可以使用余弦相似度及其对应的余弦距离：
\begin{equation}
 s_{\cos}(x,y)=\frac{x^{\mathsf T}y}{\|x\|_2\|y\|_2},\qquad
 d_{\cos}(x,y)=1-s_{\cos}(x,y),
 \quad x\ne0,\ y\ne0.\label{eq:math-cosine-distance}
\end{equation}
例如，两台设备的振动频谱一个整体强、一个整体弱，但频率峰的相对分布相似时，余弦距离可能很小，而欧氏距离仍然可能很大。相反，若要比较类别、开关状态或字符串中有多少位置不同，可以使用汉明距离：
\begin{equation}
 d_{\mathrm H}(x,y)=\sum_{j=1}^{d}\mathbb{1}\{x_j\ne y_j\},\label{eq:math-hamming-distance}
\end{equation}
它只统计不相同的位置，不表示数值差异的大小，因此不适合直接代替连续传感器变量的欧氏距离。

概率模型中的 KL 散度比较分布而非单个点：
\begin{equation}
 \KL(p\|q)=\sum_xp(x)\log\frac{p(x)}{q(x)},\label{eq:math-kl-distance}
\end{equation}
它会把概率质量放错位置的代价纳入比较，因此适合描述预测分布与真实分布之间的差异。不过，KL 散度一般不对称，也可能为无穷大。若希望比较两个分布搬运概率质量所需的“距离”，离散情形还可以使用 Wasserstein 距离：
\begin{equation}
 W_1(p,q)=\min_{\gamma\in\Gamma(p,q)}
 \sum_{a,b}\gamma(a,b)c(a,b),\label{eq:math-wasserstein-distance}
\end{equation}
其中$\Gamma(p,q)$是边缘分布分别为$p$和$q$的联合分布集合，$c(a,b)$是把质量从$a$搬到$b$的代价。举例来说，把预测的风速分布从“每秒 8 米”移到“每秒 9 米”，Wasserstein 距离会利用这两个数值相近这一事实；KL 散度则只看概率质量在各个取值上的比例。实际使用哪一种距离，取决于特征的尺度、是否关心方向、是否关心最坏坐标，以及分布取值之间是否有自然的距离。

\section{特征值、主方向与降维}
对称矩阵可以正交对角化：$A=U\Lambda U^{\mathsf T}$。若$A$是协方差矩阵，特征向量给出数据方差最大的正交方向，特征值给出对应方差。Principal Component Analysis (PCA) 选择前$r$个方向的投影
\begin{equation}
 z=U_r^{\mathsf T}(x-\bar x),
 \qquad \hat x=U_rz+\bar x.
\end{equation}
它同时可以理解为最大化投影方差，或最小化线性子空间的重构误差。这个双重解释为后文的表示学习和流形学习提供了桥梁。

\paragraph{从白化到PCA的优化证明。}
当$\Sigma$正定，令$W=\Lambda^{-1/2}U^{\mathsf T}$，则$d_M(x,y)=\|W(x-y)\|_2$。这是先把各主方向除以其标准差再测欧氏距离；但协方差同时含有信号变化和噪声变化，不能无条件把大方差方向都解释为低质量传感器。样本协方差奇异时可考虑$\Sigma+\epsilon I$，代价是改变了原距离。

对中心化样本$\bm z_i=\bm x_i-\bar{\bm x}\in\R^d$，定义$S=n^{-1}\sum_i\bm z_i\bm z_i^{\mathsf T}$。单个单位方向$\bm u$的平均投影平方为$\bm u^{\mathsf T}S\bm u$。拉格朗日函数$\bm u^{\mathsf T}S\bm u-\lambda(\bm u^{\mathsf T}\bm u-1)$的驻点满足$S\bm u=\lambda\bm u$；按特征向量展开$\bm u$可知该二次型是特征值的加权平均，最大值因此为最大特征值。

对$U_r\in\R^{d\times r}$且$U_r^{\mathsf T}U_r=I_r$，投影与残差正交，故
\begin{align}
 \frac1n\sum_i\|\bm z_i-U_rU_r^{\mathsf T}\bm z_i\|_2^2
 &=\frac1n\sum_i\bigl(\|\bm z_i\|_2^2-\|U_r^{\mathsf T}\bm z_i\|_2^2\bigr)\\
 &=\operatorname{tr}(S)-\operatorname{tr}(U_r^{\mathsf T}SU_r).\label{eq:math-pca-reconstruction}
\end{align}
式\eqref{eq:math-pca-reconstruction}中，$\operatorname{tr}$表示对角元素之和。在$S$的特征基中，后一迹是特征值按各方向保留程度的加权和，权重在$[0,1]$内且总和为$r$，所以保留最大的$r$个特征值达到最大值，最小重构误差为$\sum_{j>r}\lambda_j$。重根时最优基不唯一；大方差也不必与故障标签有关，PCA最优的是重构，不是分类。

\section{微分、积分与期望的对应}
连续概率密度的积分必须等于一；离散概率则通过求和归一化。对参数化密度$p_\theta(x)$，当分布本身也依赖参数时，采样分布的变化也必须被考虑：
\begin{equation}
 \nabla_\theta\mathbb{E}_{x\sim p_\theta}[f_\theta(x)]
 \neq \mathbb{E}_{x\sim p_\theta}[\nabla_\theta f_\theta(x)].\label{eq:math-expectation-gradient-warning}
\end{equation}
这里有两种变化要一起计算：参数改变后，同一个样本上的函数值会变，抽到各个样本的概率也可能变。普通监督学习通常把训练样本固定，只需处理前一种变化；变分推断中的采样分布也要学习，因此常用重参数化技巧或得分函数估计来处理后一种变化。

\paragraph{把缺少的梯度项补出来。}
式\eqref{eq:math-expectation-gradient-warning}的$\neq$应理解为一般不能直接交换，并非所有函数都严格不等。若积分支撑集不随$\theta$变化，密度与函数可微，且存在允许微分与积分交换的可积控制函数，则乘积法则给出
\begin{align}
 \nabla_\theta\int f_\theta(x)p_\theta(x)\,dx
 &=\int\bigl[p_\theta(x)\nabla_\theta f_\theta(x)
       +f_\theta(x)\nabla_\theta p_\theta(x)\bigr]dx\\
 &=\E_{p_\theta}[\nabla_\theta f_\theta(X)]
   +\E_{p_\theta}[f_\theta(X)\nabla_\theta\log p_\theta(X)].\label{eq:math-score-gradient}
\end{align}
式\eqref{eq:math-score-gradient}的第二项就是采样分布变化的贡献。若可写$X=\mu+\sigma\varepsilon$、$\varepsilon\sim\mathcal N(0,1)$且$\sigma>0$，也可对固定噪声分布下的$f_\theta(\mu+\sigma\varepsilon)$使用链式法则。若支撑集依赖参数，例如$X\sim\operatorname{Uniform}(0,\theta)$，则还要处理积分边界，不能机械使用式\eqref{eq:math-score-gradient}。

\section{数值稳定性}
概率模型经常使用对数概率。直接计算$\exp(a)$可能溢出，通常使用
\begin{equation}
 \log\sum_i\exp(a_i)=m+\log\sum_i\exp(a_i-m),\qquad m=\max_i a_i.
\end{equation}
这就是 log-sum-exp 技巧。softmax 的实现、交叉熵损失和扩散模型中的噪声权重都依赖类似的稳定计算。数值稳定性不是实现细节之外的问题：溢出的 loss 会让优化器得到 NaN，随后整个模型参数都可能失效。

\section{本章小结：看到公式时问什么}
面对一个新公式，先问变量是确定量还是随机量，再问每一项的维度和求和范围；然后检查它是在拟合条件均值、概率分布、几何结构还是约束；最后确认该量在训练和推理阶段是否都可计算。这个检查顺序把符号阅读转化为模型理解。

本章的几条推导可以连起来阅读：投影把逼近误差分解为可表达部分和残差；条件均值把预测误差分解为可约误差和条件波动；KL把交叉熵分解为真实分布自身的不确定性和模型失配。它们共同提醒读者，优化只能改变目标中的某些项，不能凭空补充缺失的信息，也不能消除模型假设的边界。

下一章把这些运算用到具体任务中：一段振动记录算一个样本，还是一次完整运行算一个样本？要预测平均温度、温度的某个分位数，还是故障概率？这些选择会决定标签和损失，也决定本章的公式该用在哪里。

